\section{Основы Python}
\label{sec:python_programming}

\begin{infobox}[Цели главы]
\begin{itemize}
    \item Понять, чем Python отличается от Lua, и запустить первую программу.
    \item Работать с переменными, типами, арифметикой и логикой.
    \item Управлять ходом программы: ветвления \texttt{if}/\texttt{elif}/\texttt{else}, циклы \texttt{for} и \texttt{while}.
    \item Писать функции --- строительный материал для программ полёта в следующей главе.
\end{itemize}
\end{infobox}

\begin{missionbox}[Миссии симулятора к главе]
В панели \ui{Обучение} выберите язык \ui{Python}.
\begin{itemize}
    \item \textbf{1--9 «Дрон и первая команда»}: аппарат в сцене, где живёт программа, команда автопилоту, повторный arm, остановить и начать заново, журнал первого запуска, что такое API, одна ошибка за раз, самостоятельная подготовка. Здесь впервые встречаются \verb|print| и единственная команда дрону \verb|pioneer.arm()|.
    \item \textbf{19--27 «Маленькие программы»}: имя для числа, подпись к измерению, вычисляем настройку, выбор по условию, сравниваем числа, повторяем действия, действие с именем, настройка функции, список сигналов. Это материал уроков этой главы, применённый к подсветке дрона.
\end{itemize}
В миссиях 19--27 подсветку включает команда \verb|pioneer.led_control(номер, r, g, b)|; подробно она разобрана в разделе~\ref{sec:python-leds}, а миссии 10--18 про подсветку удобнее пройти вместе со следующей главой.
\end{missionbox}

\begin{rulebox}[Как работать с главой]
Примеры запускайте в симуляторе, выбрав язык \ui{Python}: вывод \verb|print| появится в терминале, как и для Lua. Структура главы повторяет главу про Lua, поэтому сравнивайте конструкции двух языков --- это лучший способ запомнить оба. Ответы к задачам для самостоятельного решения и к заданиям уроков лежат в папке \texttt{examples/python-answers}: сначала решите сами, потом сверьтесь с файлом.
\end{rulebox}
